% A.67 reserved by the root; this statement is explicitly local.
\subsection{A noncongruence cover distinguishing equal local congruence data}

\begin{statement}{A.67}{Proved}
{An exact local point--cusp comparison in the full geometric congruence quotient}
Put $K=\Q_5$.  The point section of $5\in U(K)$ and a
cuspidal section at zero have conjugate images in the full
Legendre congruence extension $E_H$ over $G_K$, with
conjugator in the geometric group $H$.
The cuspidal section has the explicit tangent parameter
\[
                 C=16\sqrt{q_0}=5\sqrt{F(5)}\in K^\times,
 \qquad q_0=\frac{25}{256}F(5),\quad
                       \sqrt{F(5)}\in1+5\Z_5,
\tag{A.67.1}
\]
where $q_0$ is the Tate parameter of $E_5$ and $F$ is
the normalized formal series in (A.63.7).

The finite etale cover
\[
                  V\longrightarrow U_K,
                       \qquad z^5=1-t,
\tag{A.67.2}
\]
admits a lift of this cuspidal section and admits no lift
of the point section of five.  Its geometric quotient
does not factor through $H$.  Thus the two local full
congruence images are equal up to geometric conjugacy,
while this specified noncongruence neighbourhood detects
the two sections.
\end{statement}
\begin{proof}
We first specify the ordinary parameter in (A.67.1).
The inverse of $1/J(q)\in q\Z[[q]]$ has integral
coefficients, as in A.63.  Substituting the exact Legendre
formula shows more precisely that
\[
               q(t)=\frac{t^2}{256}F(t),
                      \qquad F(t)\in1+t\Z[1/2][[t]].
\tag{A.67.3}
\]
Indeed
$1/j(E_t)=t^2(1-t)^2/[256(t^2-t+1)^3]$ has the
indicated coefficients, and its leading coefficient is
$1/256$.
Consequently this series converges at $t=5$ and
$F(5)\in1+5\Z_5$.  Squaring is a bijection on
$1+5\Z_5$, by Hensel's lemma, so the square root in
(A.67.1) is uniquely defined.  It satisfies the exact
ordinary-number identity
\[
                             C^2/256=q_0.
\tag{A.67.4}
\]
For this curve
\[
 c_4=336,\qquad c_6=5184=72^2,\qquad
 \Delta=6400,\qquad j(E_5)=148176/25.
\tag{A.67.5}
\]
It has split multiplicative reduction over $K$: modulo
five the node has the two slopes $\pm\sqrt{-1}$, and
$-1$ is a square in $K$.
Thus it is the split Tate curve with parameter $q_0$
over $K$ itself.

Let $s_C^{\rm cusp}$ be the full cusp-zero section with
tangent class $C$.  Concretely, after the tangent-one
splitting of its decomposition group
$\widehat\Z(1)\rtimes G_K$, it is the continuous
section whose inertia-coordinate cocycle is $\kappa(C)$.
The cusp-zero splitting character is that of $-1$,
which is also trivial over $K$.
The compatible roots of the formal unit $F(t)$ have
constant term one, so its specialization along this
decomposition section has zero Kummer class, as proved
in A.63.  Its complete integral Tate extension is
therefore the class of $C^2/256=q_0$.

For every coefficient prime $\ell$, this gives an
isomorphism between the two integral Tate extensions
preserving their Weil pairings.
For clarity, choose the usual root-of-unity and
parameter-root Tate bases, in which
$e_N(\zeta_N^a q_0^{b/N},\zeta_N^c q_0^{d/N})
=\zeta_N^{ad-bc}$
\cite[\S19.2]{LegIovitaMasdeu2009}.
The two integral extension cocycles have the same class
in $H^1(K,\Z_\ell(1))$ by (A.67.4), so their difference
is the coboundary of an integral zero-cochain.
An upper triangular unipotent matrix for that
zero-cochain intertwines the two extensions and has
determinant one.  The curve-to-Tate isomorphisms
preserve the same pairings.
At every odd $\ell$, the resulting conjugator is
therefore an element of
$\operatorname{SL}_2(\Z_\ell)=H_\ell$.
This includes the coefficient prime five.

At two we verify the stronger geometric congruence
condition directly, so that no arbitrary permutation
of the ordered two-torsion labels enters the argument.
A.49 gives an actual geometric pro-two conjugacy
between the point section of five and the cusp-zero
section with tangent class $\kappa_2(5)$.
By (A.67.1),
\[
                         C/5\in1+5\Z_5.
\tag{A.67.6}
\]
Every power map $u\mapsto u^{2^a}$ is a bijection on
this group.  Hence $\kappa_2(C/5)=0$ and the two
cusp sections have the same geometric pro-two image
up to conjugacy.  The map from the geometric pro-two
quotient to $H_2$ is the actual Legendre monodromy map,
since $H_2$ is a pro-two group.  Thus this conjugacy
maps to one with conjugator in $H_2$.

A.17 identifies
$H=H_2\times\prod_{\ell\ {
m odd}}
\operatorname{SL}_2(\Z_\ell)$.
The product of the just constructed conjugators is
therefore in $H$.  It intertwines the full integral
Legendre representations and the two projections to
$G_K$.  The defining map
\[
                  E_H\longrightarrow
                  \operatorname{GL}_2(\widehat\Z)\times G_K
\tag{A.67.7}
\]
is injective: an element of its kernel is geometric
and has trivial Legendre image, which is precisely
the kernel divided out in defining $E_H$.
This proves the claimed conjugacy of the two actual
arithmetic sections of $E_H$.

We now use the specified cover.  Since $t\ne1$ on
$U_K$, the equation $z^5=1-t$ has $z\ne0$ and
nonzero derivative $5z^4$; it defines a finite etale
degree-five cover.
Its geometric monodromy sends the meridians about
$0,1,\infty$ to $0,1,-1$ in $\Z/5$, respectively.
It is nontrivial, as also follows from the order-one
zero of $1-t$ at one.  If this quotient factored
through $H$, it would give a nontrivial homomorphism
$H^{\rm ab}\to\Z/5$.
The explicit group
$H^{\rm ab}=(\Z/8)^2\oplus\Z/3$ in A.56 has no
such homomorphism.  Thus this cover detects an
actual part of the geometric congruence kernel.

On the cusp-zero section, the function $1-t$ is a
formal unit of value one and has compatible roots
of constant term one.  In particular its fifth
Kummer class is zero.  The sheet specializing to
$z=1$ is a $G_K$-fixed sheet of the pulled back
cover.  Its stabilizer is the corresponding open
subgroup $\pi_1(V)\subset\pi_1(U_K)$, so the full
cusp section lifts to this subgroup.

For the point section of five the pulled back
finite $G_K$-set is the fibre $z^5=-4$.
If $z\in K$ solved this equation, its valuation
would be zero, and reduction modulo five would
give $z\equiv1\pmod5$.  Writing $z=1+5a$ then
gives $z^5\equiv1\pmod{25}$, whereas
$-4\equiv21\pmod{25}$.  Hence this fibre has no
$K$-rational point and no fixed sheet, so its
point section has no lift to the cover.
This proves both lift assertions by the actual
stabilizer maps.
\end{proof}

\noindent\emph{Scope and exact arithmetic certificate.}
The section in this statement is over $\Q_5$.
No global Selmer section with these mixed local
properties is asserted.
Its tangent $C$ is transcendental over $\Q$:
if it were algebraic, (A.67.4) would make $q_0$
algebraic, contradicting the proved Mahler--Manin
theorem and the rational value (A.67.5).
The program
\path{code/local_congruence_cusp_certificate.py}
checks the exact invariants, the integral inverse
$j$ equation modulo $5^{14}$, the square-root
identity, and the fifth-power obstruction.
Its retained output has $F(5)\equiv6\pmod{25}$,
$C/5\equiv16\pmod{25}$, and the unit fifth powers
modulo twenty-five equal to $1,7,18,24$.
These finite checks accompany the full maps and
infinite identities proved above.
